\documentclass[11pt]{amsart}
\usepackage{amsmath,amssymb,amsthm}
\usepackage{microtype}
\usepackage[hidelinks]{hyperref}
\newtheorem{theorem}{Theorem}[section]
\newtheorem{lemma}[theorem]{Lemma}
\newtheorem{proposition}[theorem]{Proposition}
\newtheorem{corollary}[theorem]{Corollary}
\theoremstyle{definition}
\newtheorem{definition}[theorem]{Definition}
\newtheorem{construction}[theorem]{Construction}
\newcommand{\NE}{\mathrm{NE}}
\newcommand{\SPE}{\mathrm{SPE}}
\newcommand{\set}[1]{\{#1\}}
\newcommand{\sorted}{\operatorname{sort}}
\title[Factor-two equilibria with atomic clients]{Factor-Two Subgame-Perfect Equilibria with Atomic Clients and a Common Location Set}
\date{}
\begin{document}
\begin{abstract}
Krogmann, Lenzner, Skopalik, Uetz, and Vos introduced a two-stage facility-location game in which facilities choose locations and weighted atomic clients then choose accessible facilities to minimize congestion. Their existence guarantee for weighted clients depends on the number of facilities, while their lower bound is the golden ratio. We prove that a factor-two approximate subgame-perfect equilibrium exists whenever the facilities share a finite location set, independently of their number. Client equilibria are exact at every placement. On path, every covered client chooses an occupied location and independently randomizes uniformly among its facilities; all other continuations may be pure. We obtain the placement by lexicographically maximizing the facility-load vector. The resulting weight bounds, combined with an atomic packing argument, yield exact client equilibria after every deviation and limit every facility's improvement factor to two.
\end{abstract}
\maketitle

\section{Introduction}\label{sec:introduction}
We study approximate subgame-perfect equilibria in a facility-location game with weighted atomic clients. Facilities first choose locations, and the clients then choose accessible facilities to minimize congestion. A facility move changes both the clients it can serve and the client game. The equilibrium existence question therefore requires a choice of client equilibrium at every placement, with a bound on each facility's gain from moving.

Krogmann, Lenzner, Skopalik, Uetz, and Vos~\cite{krogmann2024} study this problem for clients with unsplittable weights who may randomize over accessible facilities. They prove existence of exact subgame-perfect equilibria for unit-weight clients. For unequal weights, their Theorem~5 gives examples excluding every approximation factor below the golden ratio, while Observation~4 gives a guarantee of $k$ for $k$ facilities. Their Section~6 asks whether the golden ratio suffices in the weighted game. We consider the case in which all facilities have the same available locations, and prove an existence guarantee of two for every $k$.

To prove the factor-two bound, we need a placement with controlled client weight after each facility move and an exact client equilibrium respecting that control. We first maximize the sorted facility-load vector over placements and a finite family of mixed client profiles. We show that the selected profile is a client equilibrium and derive bounds on the weight affected by every facility move. We then construct a pure client equilibrium after each deviation. In this construction, a client heavier than the required load bound stays alone at a stationary facility, while strict improving moves of the other clients preserve the bound. Combining these choices gives a complete continuation.

\subsection{The model and the result}\label{subsec:model}
We specify the game by its clients, locations, and coverage sets. Let $I$ be a finite client set with weights $w_i>0$, let $S$ be a finite nonempty location set, and let $F=\set{1,\ldots,k}$ be the set of labeled facilities, where $k\ge2$. A location $t\in S$ covers a specified set $C_t\subseteq I$. Every facility may choose every location in $S$. A placement is a tuple $s=(s_f)_{f\in F}\in S^k$; several facilities may occupy the same location.

At a fixed placement $s$, client $i$ may use the facilities in
\[
 A_i(s)=\set{f\in F:i\in C_{s_f}}.
\]
If $A_i(s)\ne\varnothing$, the client must choose one of these facilities; otherwise it is unserved. Clients randomize independently. A feasible client profile $p=(p_{if})$ has $p_{if}=0$ outside $A_i(s)$ and $\sum_{f\in A_i(s)}p_{if}=1$ when $A_i(s)\ne\varnothing$. For an unserved client all probabilities are zero. The expected load of facility $f$ and the conditional cost to client $i$ of using $f$ are
\begin{equation}\label{eq:cost}
 L_f(p)=\sum_{i\in I}w_ip_{if},
 \qquad c_{if}(p)=w_i+\sum_{j\ne i}w_jp_{jf}.
\end{equation}
In particular, the whole weight $w_i$ enters the cost when client $i$ uses a facility. A client profile is an exact Nash equilibrium if every facility used with positive probability minimizes the client's conditional cost among its accessible facilities.

A continuation assigns an exact client equilibrium $\sigma(t)$ to every labeled placement $t\in S^k$. We write $s[f\leftarrow r]$ for the placement obtained by moving only facility $f$ from $s_f$ to $r$.

\begin{definition}\label{def:spe}
For $\alpha\ge1$, an $\alpha$-approximate subgame-perfect equilibrium is a placement $s$ and a continuation $\sigma$ such that
\begin{equation}\label{eq:spe}
 L_f\bigl(\sigma(s[f\leftarrow r])\bigr)
 \le\alpha L_f\bigl(\sigma(s)\bigr)
 \qquad(f\in F,\ r\in S).
\end{equation}
The client equilibrium requirement is exact at every placement.
\end{definition}

We prove that every instance of this model admits a $2$-approximate subgame-perfect equilibrium. On path, every covered client chooses a single occupied location and independently samples a facility there uniformly. At every other placement we may choose a pure client equilibrium. The proof gives a finite exact construction; a polynomial-time algorithm for general inputs remains open.

The common location set is used when a facility is moved to another occupied location to compare two load vectors. Thus the result concerns a restriction of the facility-dependent location sets allowed in~\cite{krogmann2024}. Their two-facility lower bound still leaves a gap between the golden ratio and two for a uniform guarantee over all numbers of facilities.

\noindent\textbf{Organization.}
Section~\ref{sec:proof} derives the weight bounds from a lexicographic maximum, converts those bounds into exact client equilibria after deviations, and assembles the complete continuation. It concludes with the finite certificate supplied for rational input.

\noindent\textbf{Conventions and notation.}
For $J\subseteq I$, write $w(J)=\sum_{i\in J}w_i$. We sort load vectors in nondecreasing order and compare the resulting vectors lexicographically. Facility labels are retained in every placement and continuation. The weights are positive real numbers except in the explicit encoding statement at the end of Section~\ref{sec:proof}.

\section{A factor-two equilibrium}\label{sec:proof}
To obtain a complete continuation, we must bound the moving facility's load at an exact client equilibrium after every deviation. We first derive sufficient weight bounds from a lexicographically maximal placement and assignment. An atomic packing lemma then gives a pure client profile after each move. We show that strict improving client moves preserve the required load bound and terminate at an exact equilibrium.

\subsection{Lexicographic selection and transfer bounds}\label{subsec:selection}
At an occupied location, independent uniform randomization gives equal expected facility loads while keeping each client's weight indivisible in every realization. We maximize the resulting load vector over a finite family of profiles.

\begin{construction}\label{con:selection}
For any placement $s$, let $O=\set{t\in S:s_f=t\text{ for some }f}$ and $q_t=|\set{f:s_f=t}|$ for $t\in O$. Assign every covered client $i$ to a single location $t(i)\in O$ that covers it, and put
\[
 J_t=\set{i:t(i)=t},\qquad W_t=w(J_t).
\]
Each client in $J_t$ chooses each of the $q_t$ facilities at $t$ with probability $1/q_t$, independently of all other clients. Every facility at $t$ then has expected load $W_t/q_t$.

Fix a total order on the finite encodings of placements and site assignments. Choose the first pair in this order for which the sorted vector of all $k$ facility loads is lexicographically maximal over all placements and all such assignments. Denote its client profile by $p$.
\end{construction}

The maximum exists because there are finitely many placements and site assignments. The following elementary comparison will be used repeatedly: replacing one or more occurrences of a value $b$ by values strictly above $b$, while all changed entries previously above $b$ remain above $b$ and all remaining entries do not decrease, strictly increases the sorted vector lexicographically. The comparison remains valid if some smaller entries increase, since their first increase already gives a strict improvement.

\begin{lemma}\label{lem:onpath}
If $R=\max_{t\in S}w(C_t)>0$, the profile $p$ in Construction~\ref{con:selection} is an exact client equilibrium and every facility has load at least $R/k$.
\end{lemma}
\begin{proof}
Placing all facilities at a location of reach $R$ and using the prescribed uniform probabilities gives each facility load $R/k$. Lexicographic maximality therefore implies that the smallest load in the selected profile is at least $R/k$.

Consider a client $i\in J_t$. Its conditional cost at any facility at $t$ is
\begin{equation}\label{eq:sitecost}
 w_i+\frac{W_t-w_i}{q_t}.
\end{equation}
At a facility at another accessible occupied location $v$, its cost is $w_i+W_v/q_v$. Suppose this alternative were strictly better, so that
\[
 \frac{W_t-w_i}{q_t}>\frac{W_v}{q_v}=b.
\]
Reassign only client $i$ from $t$ to $v$ in Construction~\ref{con:selection}. All $q_v$ coordinates formerly equal to $b$ become $(W_v+w_i)/q_v>b$, while the $q_t$ coordinates at $t$ remain above $b$. This strictly increases the sorted vector, a contradiction. Facilities at $t$ have equal conditional costs by~\eqref{eq:sitecost}. Hence no pure client deviation improves, and linearity of expected cost in the client's own probabilities excludes a mixed improvement as well.
\end{proof}

The selected assignment also controls clients who were unserved before a move and clients whose assigned location disappears. For an unoccupied location $r\notin O$, define
\begin{equation}\label{eq:newweight}
 N_r=w\left(C_r\setminus\bigcup_{t\in O}C_t\right).
\end{equation}
These are precisely the clients newly covered when a facility moves to $r$.

\begin{lemma}\label{lem:transfer}
Assume $R>0$. Fix a facility $f$ at $u=s_f$, and put $a=W_u/q_u$. If $q_u\ge2$, then
\begin{align}
 W_t&\le(q_t+1)a &&(t\in O\setminus\set{u}),\label{eq:occupied-survive}\\
 N_r&\le a &&(r\notin O).\label{eq:empty-survive}
\end{align}
If $q_u=1$, put $J=J_u$. Then
\begin{align}
 W_t+w(J\cap C_t)&\le(q_t+1)a &&(t\in O\setminus\set{u}),\label{eq:occupied-disappear}\\
 N_r+w(J\cap C_r)&\le a &&(r\notin O).\label{eq:empty-disappear}
\end{align}
\end{lemma}
\begin{proof}
 Suppose first that $q_u\ge2$. Moving $f$ to $t\in O\setminus\set{u}$ and keeping all client site assignments feasible gives the remaining facilities at $u$ load $q_ua/(q_u-1)>a$. If~\eqref{eq:occupied-survive} failed, all $q_t+1$ facilities at $t$ would also have load above $a$. The old target load was above $a$ as well, since $W_t/q_t>((q_t+1)a)/q_t>a$. Every coordinate equal to $a$ at $u$ would thus be replaced by a larger one, with no smaller coordinate decreasing. This contradicts maximality. If~\eqref{eq:empty-survive} failed, moving $f$ to $r$ and assigning all newly covered clients to it would replace its load $a$ by $N_r>a$ and increase the remaining loads at $u$, again a contradiction.

Now suppose that $q_u=1$, so $w(J)=a$. If~\eqref{eq:occupied-disappear} failed, move $f$ to $t$ and assign all clients of $J\cap C_t$ to $t$. Clients in $J\setminus C_t$ that remain covered are assigned to any other remaining occupied location that covers them; those that become uncovered are unserved. Clients outside $J$ retain their assignments. The new load at $t$ is above $a$. Its old load is also above $a$, since
\[
 W_t+w(J\cap C_t)>(q_t+1)a,
 \quad w(J\cap C_t)\le a
 \quad\Longrightarrow\quad W_t>q_ta.
\]
The coordinate $a$ at $u$ disappears, every changed coordinate at $t$ remains above $a$, and no other facility load decreases. This strictly improves the sorted vector.

Finally, if~\eqref{eq:empty-disappear} failed, move $f$ to $r$ and assign both the newly covered clients and $J\cap C_r$ to $r$. These two sets are disjoint. Redistribute the other still-covered clients of $J$ to any remaining locations that cover them. The coordinate $a$ is replaced by $N_r+w(J\cap C_r)>a$, while all other coordinates do not decrease. This is the final contradiction.
\end{proof}

The terms involving $J$ account for the clients assigned to a location with only one facility. Including them before choosing a destination permits an arbitrary redistribution of those clients after that facility departs.

\subsection{Packing clients and equilibrating after a move}\label{subsec:repair}
We now distribute the weight bounded in Lemma~\ref{lem:transfer} over the stationary facilities. The next lemma allows a client heavier than the desired cap to occupy a facility alone. All remaining facilities satisfy the cap.

\begin{lemma}\label{lem:packing}
Let $a>0$ and let $h\ge1$ be an integer. Any finite list of positive weights with total at most $(h+1)a$ can be partitioned into at most $h$ parts, each having weight at most $2a$ or consisting of a single client of weight greater than $2a$.
\end{lemma}
\begin{proof}
Start with one part per client and merge two parts whenever their combined weight is at most $2a$. A part heavier than $2a$ remains a singleton. If the process stopped with $b\ge h+1$ parts, every pair of distinct parts would have combined weight above $2a$. Writing their weights as $B_1,\ldots,B_b$ and summing over pairs gives
\[
 (b-1)\sum_{j=1}^bB_j
 >2a\binom b2=(b-1)ba.
\]
Consequently the total weight would exceed $ba\ge(h+1)a$, a contradiction. The empty list gives no parts.
\end{proof}

The assignment obtained by packing need not be a client equilibrium. We next show that strict improving client moves preserve the cap on every facility that does not host an isolated heavy client. Here loads refer to a pure client assignment.

\begin{lemma}\label{lem:capped}
Fix a placement, a facility $f$, and a number $a>0$. Suppose there is a feasible pure client assignment in which $f$ has load at most $2a$, every client heavier than $2a$ is alone at a facility other than $f$, and every facility that does not host one of these heavy clients has load at most $2a$. Then this placement has a pure exact client equilibrium in which $f$ has load at most $2a$.
\end{lemma}
\begin{proof}
Repeatedly perform a strict improving pure client move. An isolated client of weight $w_i>2a$ already incurs its smallest possible cost $w_i$, so it never moves strictly profitably. A client at a facility of load at most $2a$ cannot improve by joining one of these isolated clients, because its destination cost would exceed $2a$. If it moves from facility $h$ to facility $g$, strict improvement means
\[
 L_g+w_i<L_h\le2a.
\]
Thus its destination and source loads remain at most $2a$, and the hypotheses are invariant.

For the potential $\Psi=\tfrac12\sum_{v\in F}L_v^2$, a move by $i$ from $h$ to $g$ changes its value by
\begin{equation}\label{eq:potential}
 \Psi'-\Psi=w_i(L_g+w_i-L_h)<0.
\end{equation}
There are finitely many feasible pure assignments, so the process terminates. No client then has a profitable pure move; since the other clients are fixed, randomization cannot improve its cost either. The terminal assignment is the required exact client equilibrium.
\end{proof}

Lemmas~\ref{lem:packing} and~\ref{lem:capped} now apply to every deviation of a facility in the lexicographically selected placement.

\begin{lemma}\label{lem:deviation}
Assume $R>0$. For every facility $f$ and every location $r\ne s_f$, the placement $s[f\leftarrow r]$ has an exact pure client equilibrium $p'$ satisfying
\begin{equation}\label{eq:bound-deviation}
 L_f(p')\le2L_f(p).
\end{equation}
\end{lemma}
\begin{proof}
Put $u=s_f$ and $a=L_f(p)=W_u/q_u$. Facilities other than $f$ are called stationary. We construct a feasible pure assignment satisfying Lemma~\ref{lem:capped}.

\textbf{Case 1: $q_u\ge2$.}
Every old occupied location survives. Keep each previously served client at its assigned location. There are $h_u=q_u-1$ stationary facilities at $u$, and the assigned weight there is
\[
 W_u=q_ua=(h_u+1)a.
\]
At each other old location $t$, there are $h_t=q_t$ stationary facilities, and~\eqref{eq:occupied-survive} gives $W_t\le(h_t+1)a$. Apply Lemma~\ref{lem:packing} at each location and assign the parts to its stationary facilities, leaving unused facilities empty. Every client is eligible for every facility at its assigned location.

If $r$ was unoccupied, assign all newly covered clients to $f$; their total weight is at most $a$ by~\eqref{eq:empty-survive}. If $r$ was occupied, leave $f$ empty. This serves every covered client. Every facility other than $f$ either has load at most $2a$ or hosts a single client heavier than $2a$.

\textbf{Case 2: $q_u=1$.}
The old location $u$ disappears. Put $J=J_u$ and keep clients outside $J$ at their old assigned locations. Assign each client of $J$ that has a stationary accessible facility to any such location. If no stationary facility is accessible but $r$ covers the client, assign it to $f$. All remaining clients of $J$ become unserved.

At a surviving old location $t$, the added clients form a subset of $J\cap C_t$. Together with the original clients $J_t$, their total weight is at most $(q_t+1)a$ by~\eqref{eq:occupied-disappear}. Apply Lemma~\ref{lem:packing} to all clients now assigned to $t$ and place the resulting parts on its $q_t$ stationary facilities.

If $r$ was unoccupied, assign the genuinely newly covered clients to $f$ as well. Its load is then at most $N_r+w(J\cap C_r)\le a$ by~\eqref{eq:empty-disappear}. If $r$ was already occupied, every client of $J$ that remains covered has a stationary option, so $f$ receives no client initially. Once again every covered client is served, every client heavier than $2a$ is isolated at a stationary facility, and every non-heavy facility has load at most $2a$.

In both cases Lemma~\ref{lem:capped} supplies an exact pure client equilibrium with $L_f\le2a$, proving~\eqref{eq:bound-deviation}.
\end{proof}

\subsection{Completion of the continuation}\label{subsec:completion}
It remains to specify one continuation at all placements. Since distinct unilateral deviations from a fixed labeled placement reach distinct placements, the client equilibria supplied by Lemma~\ref{lem:deviation} can be selected simultaneously.

\begin{theorem}\label{thm:main}
For every finite nonempty common location set, every finite set of clients with positive real weights, and every integer $k\ge2$, there exists a $2$-approximate subgame-perfect equilibrium. On path every covered client may choose one occupied location and independently randomize uniformly among its facilities. Every off-path client equilibrium may be chosen pure.
\end{theorem}
\begin{proof}
If $R=0$, all facility loads are zero at every placement and the assertion is immediate. Otherwise choose $(s,p)$ as in Construction~\ref{con:selection}. Lemma~\ref{lem:onpath} gives an exact on-path client equilibrium, and Lemma~\ref{lem:deviation} gives a pure client equilibrium satisfying~\eqref{eq:spe} with $\alpha=2$ at each actual deviation. These placements are pairwise distinct: deviations by different facility labels change different coordinates of $s$, and different targets for one facility change its coordinate to different values. At every remaining placement, start with any feasible pure assignment and perform strict improving client moves. Equation~\eqref{eq:potential} ensures termination at a pure client equilibrium. These choices define the required complete continuation.
\end{proof}

The on-path construction gives equal expected loads to facilities at the same location. After a deviation we use a different selection: an isolated heavy client remains at a stationary facility, and the moving facility satisfies the factor-two load bound. The continuation need not preserve symmetry among facilities at a location.

For rational input, we can record the selected equilibria with polynomially many bits and use a fixed rule to complete all other placements. The next proposition states this encoding consequence.

\begin{proposition}\label{prop:certificate}
Suppose the weights are positive binary rationals and the client--location incidences and all $k$ facilities are explicit. There is a finite exact procedure that outputs the placement in Theorem~\ref{thm:main}, its on-path client equilibrium, a pure client equilibrium for each of its $k(|S|-1)$ actual deviations, and a deterministic default-rule program for all other placements. The displayed strategies have encoding length polynomial in the input length. The stored on-path and actual-deviation profiles, and the facility inequalities for those deviations, can be checked in time polynomial in their explicit encoding length; correctness of the default-rule program follows from the potential argument.
\end{proposition}
\begin{proof}
Enumerate all placements and all feasible client site assignments in the fixed order from Construction~\ref{con:selection}, and retain the first lexicographic maximum. The on-path probabilities belong to $\set{0}\cup\set{1/q:1\le q\le k}$. At each actual deviation, perform the mergers of Lemma~\ref{lem:packing} and strict improving client moves as in Lemma~\ref{lem:capped}. These processes terminate and their output probabilities are zero or one.

For the default rule, fix an enumeration of the clients. Assign each covered client to its least-index accessible facility and repeatedly choose the first client in this enumeration with a strict improvement, moving it to its least-index cost-minimizing accessible facility. Equation~\eqref{eq:potential} proves that this rule terminates at every placement. Store its fixed description together with the displayed on-path and actual-deviation profiles. Each displayed pure profile needs at most one facility index per client; the number of displayed profiles is polynomial in the explicit input length. The rational loads are sums of input weights, possibly divided by an integer at most $k$. A denominator for any one load is obtained by multiplying the product of the input denominators by an integer at most $k$. Its bit length is therefore at most the sum of the input denominator bit lengths plus $O(\log k)$, so exact comparisons also have polynomial bit length. Formula~\eqref{eq:cost} checks each displayed client equilibrium, and~\eqref{eq:spe} checks the facility deviations.
\end{proof}

Proposition~\ref{prop:certificate} distinguishes a short description of the continuation from an efficient search procedure. Exhaustive selection and finite improvement suffice for existence; neither step has a polynomial running-time bound in this proof. Determining the best uniform approximation factor between the golden ratio and two, and finding an efficient construction for general common location sets, remain separate questions.

\enlargethispage{3\baselineskip}
\bibliographystyle{plain}
\bibliography{references}
\end{document}
